% Shared body of the three competitor setup manuals. The including file defines
% \OSname, \OSkey, \OSappdir, \OSterm and \OSdockerfix, then inputs this.

\section*{What you are setting up}

A simulator of the RoboRacer 1:10 car runs natively on your \OSname. Your code runs in the box: a
Docker container with ROS~2, the bridge to the simulator and the real car's drivers, which keeps
what you install and sees your FEB folder. You never install ROS on the \OSname{} itself.
Foxglove shows you the car's sensors live. \texttt{docs/devkit.md} is the one-page reference.

\begin{lstlisting}
~/FEB/                      your folder; the box sees it at /home/feb/FEB
  feb-racing/               the repo, your clone
    feb                     the launcher
    stack/                  YOUR code (gitignored here: keep it in a repo of your own)
    starter_kit/feb_driver/ the template you copy into stack/
    tracks/                 track folders
~/.feb-sim/app/<os>/        the simulator app, downloaded by setup
Docker                      runs the box
\end{lstlisting}

The loop: \texttt{feb up} builds \texttt{stack/} in the box and starts the house nodes; they turn
the simulator's stream into the car's topics, \texttt{/scan}, \texttt{/odom}, \texttt{/imu}; your node
reads them and publishes \texttt{/drive}, a speed and a steering angle; the supervisor passes it to
the adapter, the adapter to the simulator; the car moves. On the real car the same node publishes
the same topic and the VESC does the rest.

\input{setup_install_\OSkey}

\section{Drive manually}

\begin{lstlisting}
./feb up
\end{lstlisting}

Arrow keys drive. Menu top-left: Manual/Autonomous, Track, Cars, Look, camera. Close the
window to stop. Nothing drives by itself without a stack.

\section{Drive autonomously with the starter kit}

\begin{lstlisting}
cp -r starter_kit/feb_driver stack/
./feb up --stack "ros2 launch feb_driver drive.launch.py"
\end{lstlisting}

The box builds the package, the simulator opens, the car drives, Foxglove opens on its topics.

\textbf{What the starter driver is.} A complete ROS~2 package. The driver is one commented file,
\texttt{stack/feb\_driver/feb\_driver/reactive\_driver.py}, a ``follow the gap'' controller: clean
the lidar scan, blank a safety bubble around the nearest obstacle, aim at the widest gap, pick a
speed from the room ahead, publish the speed and steering angle on \texttt{/drive}. Its numbers are in
\texttt{stack/feb\_driver/config/driver.yaml}.

\textbf{Editing code.} \texttt{stack/} is a plain folder on your \OSname; open it in any editor, or
\texttt{./feb code} opens VS~Code attached to the box, where every terminal is already inside it.
Edit, save, run the same \texttt{./feb up} command again; the box rebuilds your packages in
seconds, and a Python edit needs no rebuild at all.

\section{Score and post}

\begin{lstlisting}
./feb practice --stack "ros2 launch feb_driver drive.launch.py"
./feb submit
\end{lstlisting}

\texttt{practice} runs one attempt the way a competition does: headless simulator, a warm-up lap
plus ten timed laps, 10~s per collision, 300~s limit; it writes \texttt{runs/<id>/result.json} and a
bag. \texttt{submit} opens a pull request with that file; a check validates and merges it, and your
time appears on the site's practice board as \emph{unverified}. The organiser's nightly rerun of your
registered image on the same track turns it \emph{verified}.

\section{Looking inside: extra terminals and Foxglove}

\subsection*{More terminals in the box}

The box is always there. Open as many \OSterm{} windows as you like and, in each, from the repo
folder:

\begin{lstlisting}
./feb shell
\end{lstlisting}

That is a ROS~2 shell inside the box, with the environment already sourced. Useful commands there:

\begin{lstlisting}
ros2 topic list                       # everything on the graph
ros2 topic hz /scan                   # is the sensor stream alive, at what rate
ros2 topic echo /drive                # what your node is sending
ros2 topic echo /feb/state            # what the supervisor is doing with it
ros2 node list                        # your node should be here
ros2 param list /reactive_driver      # live parameters
ros2 bag record -a -o ~/FEB/bags/mybag   # lands on your laptop
\end{lstlisting}

\texttt{exit} leaves the shell; the box keeps running.

\textbf{The pieces on their own.} \texttt{feb up} starts the house nodes, the simulator and
Foxglove together because that is what you want most of the time, but they are independent.
The house nodes alone, for replaying bags or tests: \texttt{./feb up --no-sim}, or in a box
terminal \texttt{ros2 launch feb\_tools up.launch.py}. The simulator alone is the app in
\texttt{\textasciitilde/.feb-sim/app/\OSappdir}, keyboard driving, no ROS; it connects to the
box by itself whenever the house nodes are up.

\subsection*{Foxglove}

Installed in the install list; sign in once. After that, \texttt{./feb up} opens
Foxglove connected to your car a few seconds after the simulator window. If it does not, or
you want a second window: in Foxglove choose \emph{Open connection}, \emph{Foxglove WebSocket},
address \texttt{ws://localhost:8765}.

Panels worth adding (the \emph{+} at the top right of a panel):
\begin{itemize}
  \item \textbf{3D}: add the topic \texttt{/scan}; display frame \texttt{laser}. The scan appears as points around the car.
  \item \textbf{Image}: \texttt{/camera/color/image\_raw}.
  \item \textbf{Plot}: \texttt{/drive.drive.speed} and \texttt{/odom.twist.twist.linear.x} against time: what you asked for and what you got.
  \item \textbf{Raw Messages}: any topic, to read exact values.
  \item \textbf{Topics} (left sidebar): every topic with its rate; a dash means nothing is publishing.
\end{itemize}

Save the layout (\emph{Layout}, \emph{Save}) and it comes back next time. Foxglove also opens the
bags that \texttt{practice} records under \texttt{runs/<id>/}, and any bag you record yourself:
\emph{Open local file}.

\section{When something goes wrong}

\begin{lstlisting}
./feb logs -n 100         # build errors, your node's output
./feb shell               # a ROS 2 terminal inside the box
./feb reset               # a fresh box; your files are untouched
\end{lstlisting}

Nothing drives: correct without \texttt{--stack}; with one, read \texttt{./feb logs}.
Port 4567 in use: \texttt{./feb reset}. Docker pull denied, or ``cannot connect to the Docker daemon'': \OSdockerfix.

\section{Keeping up to date}

\begin{lstlisting}
./feb update
\end{lstlisting}

One command pulls the repo (tracks, launcher, docs, new starter packages), pulls a newer image
if there is one and replaces the box with it (what you had installed by hand in the box goes,
your files do not), and replaces the simulator app when a new build is released. Your
\texttt{stack/} folder is never touched. When a new starter package appears in
\texttt{starter\_kit/}, copy it the same way as the first one. The cone tracks
(\texttt{loop\_cones}, \texttt{spielberg\_cones}) have their own baseline, the cone driver:

\begin{lstlisting}
cp -r starter_kit/feb_cone_driver stack/
./feb up --track loop_cones --look visual --stack "ros2 launch feb_cone_driver drive.launch.py"
\end{lstlisting}

It needs the Visual look: the camera's colour thresholds are tuned for the dressed scene.
\texttt{docs/starter-kit.md} describes what it does. Whatever stack you run, \cmd{./feb logs -f}
in a second terminal shows what \texttt{feb up} printed: the package build, the house nodes, and every
line your nodes print.

\section{What the car gives and takes}

\begin{tabularx}{\linewidth}{lX}
\toprule
Topic & Meaning \\
\midrule
\texttt{/scan} & \texttt{LaserScan}, 1080 beams over 270\textdegree, 0.06 to 10~m, frame \texttt{laser} \\
\texttt{/odom} & \texttt{Odometry}: pose and speed dead-reckoned from the wheels and the gyro, as the car's VESC does it \\
\texttt{/imu} & orientation, angular velocity, linear acceleration \\
\texttt{/camera/color/image\_raw} & the front camera \\
\texttt{/drive} & \textbf{you publish this}: \texttt{AckermannDriveStamped}, \texttt{drive.speed} in m/s, \texttt{drive.steering\_angle} in rad, at most $\pm 0.5236$ \\
\texttt{/feb/state} & the supervisor: IDLE, MANUAL, AUTONOMOUS or STOP, and why \\
\texttt{/autodrive/roboracer\_1/*} & the simulator's own topics: ground truth, lap and collision counters. Debugging only, barred at race time; the car does not have them \\
\bottomrule
\end{tabularx}

\smallskip
Publish as many topics of your own as you like; only \texttt{/drive} reaches the car. Use
message stamps or \texttt{use\_sim\_time}, never the wall clock, for speeds and integrals. A
command older than half a second stops the car.

\appendix
% Appendix shared by the competitor manuals: the starter package, file by file, and the
% driver's logic per scan. Needs tikz and listings loaded by the including document.

\section{The starter package, file by file}
\label{app:starter}

\texttt{starter\_kit/feb\_driver} is a normal ROS~2 Python package. After
\texttt{cp -r starter\_kit/feb\_driver stack/} these are the files you own (the Python files live in the inner \texttt{feb\_driver/} folder):

\begin{tabularx}{\linewidth}{p{4.4cm}X}
\toprule
File & What it is \\
\midrule
\texttt{package.xml} & The ROS~2 manifest: the package name \texttt{feb\_driver}, its version and
its dependencies (\texttt{rclpy}, \texttt{sensor\_msgs}, \texttt{std\_msgs}). \texttt{colcon}
reads it to know the package exists. You only touch it to add a dependency. \\
\texttt{setup.py} & Python packaging. The line that matters is the entry point
\texttt{reactive\_driver = feb\_driver.reactive\_driver:main}: it is what makes
\texttt{ros2 run feb\_driver reactive\_driver} exist. It also installs \texttt{launch/} and
\texttt{config/} so launch files can find them. Add a line here for every new node. \\
\texttt{setup.cfg} & Two lines telling \texttt{colcon} where to put the installed scripts. Never edited. \\
\texttt{resource/feb\_driver} & An empty marker file ROS~2 uses to index the package. Never edited. \\
\texttt{\_\_init\_\_.py} & Empty; makes the folder a Python module. \\
\texttt{reactive\_driver.py} & The driver, about 170 lines, explained below. \\
\texttt{config/driver.yaml} & Every tuning number, one per line with a comment. The launch file
hands it to the node, so you change behaviour here without touching code. \\
\texttt{launch/drive.launch.py} & Starts one node, \texttt{reactive\_driver}, with the yaml;
\texttt{use\_sim\_time} follows the box (simulator clock on a laptop, wall clock on the car). This
is what \texttt{--driver feb\_driver} runs. \\
\texttt{driver.yaml} & The driver's card: its name, the kind of track it drives, its modes
(\texttt{trackdrive}, \texttt{inspection}), each a launch file and a speed cap. \\
\texttt{sysid\_probe.py}, \texttt{launch/sysid.launch.py} & Stage 2. Drives with the
same steering but a fixed throttle schedule (three steps, a coast, a chirp) while recording a
bag, to measure how the car responds. \\
\texttt{tools/sysid\_fit.py} & Fits that bag: steady-state gain (the adapter's
\texttt{speed\_per\_throttle}), idle brake, throttle lag and effective wheelbase. Run in the box. \\
\bottomrule
\end{tabularx}

\medskip
The node subscribes to two topics and publishes one, the real car's: in, \texttt{/scan}
(\texttt{LaserScan}, 1081 ranges over $\pm135^\circ$) and \texttt{/odom} (\texttt{Odometry}, the
car's own speed); out, \texttt{/drive} (\texttt{AckermannDriveStamped}): a speed in m/s and a
steering angle in rad, at most $\pm 0.5236$ ($30^\circ$). Turning the speed into throttle is the
sim adapter's job in the simulator and the VESC's on the car.

\section{How the driver drives, one scan at a time}
\label{app:driver}

Every lidar scan, about ten times a second, \texttt{on\_scan} runs the steps below and
publishes one speed and one steering angle. The figure shows one scan from above, in a corridor that bends left; the right wall is the nearest thing the lidar sees.

\begin{center}
\begin{tikzpicture}[scale=1.1, every node/.style={font=\small}]
  \begin{scope}
  \clip (-4.0,-1.5) rectangle (4.4,5.0);
  \draw[line width=1.8pt, gray!75] (1.000,-2.800) -- (1.000,2.200) -- (1.000,2.200) -- (0.994,2.370) -- (0.978,2.539) -- (0.950,2.707) -- (0.911,2.873) -- (0.862,3.036) -- (0.802,3.195) -- (0.732,3.350) -- (0.652,3.500) -- (0.562,3.644) -- (0.463,3.783) -- (0.355,3.914) -- (0.238,4.038) -- (0.114,4.155) -- (-0.017,4.263) -- (-0.156,4.362) -- (-0.300,4.452) -- (-0.450,4.532) -- (-0.605,4.602) -- (-0.764,4.662) -- (-0.927,4.711) -- (-1.093,4.750) -- (-1.261,4.778) -- (-1.430,4.794) -- (-1.600,4.800);
  \draw[line width=1.8pt, gray!75] (-1.600,-2.800) -- (-1.600,1.400) -- (-1.600,1.400) -- (-1.602,1.452) -- (-1.607,1.504) -- (-1.615,1.556) -- (-1.627,1.607) -- (-1.642,1.657) -- (-1.661,1.706) -- (-1.683,1.754) -- (-1.707,1.800) -- (-1.735,1.844) -- (-1.765,1.887) -- (-1.799,1.927) -- (-1.834,1.966) -- (-1.873,2.001) -- (-1.913,2.035) -- (-1.956,2.065) -- (-2.000,2.093) -- (-2.046,2.117) -- (-2.094,2.139) -- (-2.143,2.158) -- (-2.193,2.173) -- (-2.244,2.185) -- (-2.296,2.193) -- (-2.348,2.198) -- (-2.400,2.200);
  \fill[febblue, opacity=0.13] (0,0) -- (-1.600,0.000) -- (-1.600,0.140) -- (-1.600,0.282) -- (-1.600,0.429) -- (-1.600,0.582) -- (-1.600,0.746) -- (-1.600,0.924) -- (-1.600,1.120) -- (-1.600,1.343) -- (-1.636,1.636) -- (-3.857,4.596) -- (-3.441,4.915) -- (-3.000,5.196) -- (-2.536,5.438) -- (-2.052,5.638) -- (-1.281,4.780) -- (-0.825,4.681) -- (-0.394,4.502) -- (0.000,4.249) -- (0.344,3.926) -- (0.625,3.543) -- (0.834,3.111) -- (0.961,2.640) -- (1.000,2.145) -- (1.000,1.732) -- (1.000,1.428) -- (1.000,1.192) -- (1.000,1.000) -- (1.000,0.839) -- (1.000,0.700) -- (1.000,0.577) -- (1.000,0.466) -- cycle;
  \draw[gray!30, very thin] (0,0) -- (-1.600,-1.600);
  \draw[gray!30, very thin] (0,0) -- (-1.600,-1.343);
  \draw[gray!30, very thin] (0,0) -- (-1.600,-1.120);
  \draw[gray!30, very thin] (0,0) -- (-1.600,-0.924);
  \draw[gray!30, very thin] (0,0) -- (-1.600,-0.746);
  \draw[gray!30, very thin] (0,0) -- (-1.600,-0.582);
  \draw[gray!30, very thin] (0,0) -- (-1.600,-0.429);
  \draw[gray!30, very thin] (0,0) -- (-1.600,-0.282);
  \draw[gray!30, very thin] (0,0) -- (-1.600,-0.140);
  \draw[gray!55, thin] (0,0) -- (-1.600,0.000);
  \draw[gray!55, thin] (0,0) -- (-1.600,0.140);
  \draw[gray!55, thin] (0,0) -- (-1.600,0.282);
  \draw[gray!55, thin] (0,0) -- (-1.600,0.429);
  \draw[gray!55, thin] (0,0) -- (-1.600,0.582);
  \draw[gray!55, thin] (0,0) -- (-1.600,0.746);
  \draw[gray!55, thin] (0,0) -- (-1.600,0.924);
  \draw[gray!55, thin] (0,0) -- (-1.600,1.120);
  \draw[gray!55, thin] (0,0) -- (-1.600,1.343);
  \draw[gray!55, thin] (0,0) -- (-1.636,1.636);
  \draw[gray!55, thin] (0,0) -- (-3.857,4.596);
  \draw[gray!55, thin] (0,0) -- (-3.441,4.915);
  \draw[gray!55, thin] (0,0) -- (-3.000,5.196);
  \draw[gray!55, thin] (0,0) -- (-2.536,5.438);
  \draw[gray!55, thin] (0,0) -- (-2.052,5.638);
  \draw[gray!55, thin] (0,0) -- (-1.281,4.780);
  \draw[gray!55, thin] (0,0) -- (-0.825,4.681);
  \draw[gray!55, thin] (0,0) -- (-0.394,4.502);
  \draw[gray!55, thin] (0,0) -- (0.000,4.249);
  \draw[gray!55, thin] (0,0) -- (0.344,3.926);
  \draw[gray!55, thin] (0,0) -- (0.625,3.543);
  \draw[gray!55, thin] (0,0) -- (0.834,3.111);
  \draw[gray!55, thin] (0,0) -- (0.961,2.640);
  \draw[gray!55, thin] (0,0) -- (1.000,2.145);
  \draw[gray!55, thin] (0,0) -- (1.000,1.732);
  \draw[gray!55, thin] (0,0) -- (1.000,1.428);
  \draw[gray!55, thin] (0,0) -- (1.000,1.192);
  \draw[gray!55, thin] (0,0) -- (1.000,1.000);
  \draw[gray!55, thin] (0,0) -- (1.000,0.839);
  \draw[gray!55, thin] (0,0) -- (1.000,0.700);
  \draw[gray!55, thin] (0,0) -- (1.000,0.577);
  \draw[gray!55, thin] (0,0) -- (1.000,0.466);
  \draw[gray!30, very thin] (0,0) -- (1.000,-0.087);
  \draw[gray!30, very thin] (0,0) -- (1.000,-0.176);
  \draw[gray!30, very thin] (0,0) -- (1.000,-0.268);
  \draw[gray!30, very thin] (0,0) -- (1.000,-0.364);
  \draw[gray!30, very thin] (0,0) -- (1.000,-0.466);
  \draw[gray!30, very thin] (0,0) -- (1.000,-0.577);
  \draw[gray!30, very thin] (0,0) -- (1.000,-0.700);
  \draw[gray!30, very thin] (0,0) -- (1.000,-0.839);
  \draw[gray!30, very thin] (0,0) -- (1.000,-1.000);
  \draw[red!70, thick] (0,0) -- (1.000,0.364);
  \draw[red!70, thick] (0,0) -- (1.000,0.268);
  \draw[red!70, thick] (0,0) -- (1.000,0.176);
  \draw[red!70, thick] (0,0) -- (1.000,0.087);
  \draw[red!70, thick] (0,0) -- (1.000,0.000);
  \fill[red] (1.000,0.000) circle (2.4pt);
  \draw[red, dashed, thick] (1.000,0.000) circle (0.400);
  \draw[febblue, line width=1.6pt] (0,0) -- (-3.857,4.596);
  \draw[febblue!80, thick, dash pattern=on 3pt off 2pt] (0,0) -- (-0.825,4.681);
  \draw[-{Stealth[length=8pt]}, line width=2pt, black] (0,0) -- (-1.080,2.031);
  \draw[green!45!black, dashed] (-0.235,0) -- (-0.235,4.15); \draw[green!45!black, dashed] (0.235,0) -- (0.235,4.15);
  \fill[black!80] (-0.135,-0.27) rectangle (0.135,0.27);
  \end{scope}
  \node[anchor=west] (n) at (2.90,-0.20) {nearest return, 1.0 m};
  \draw[gray!60, thin] (n.west) -- (1.000,0.000);
  \node[anchor=west] (n) at (2.90,-0.60) {bubble 0.4 m: beams within $\pm22^\circ$ blanked};
  \draw[gray!60, thin] (n.west) -- (0.700,0.000);
  \node[anchor=east] (n) at (-3.60,1.00) {widest gap};
  \draw[gray!60, thin] (n.east) -- (-0.960,0.554);
  \node[anchor=east] (n) at (-3.60,3.90) {deepest beam};
  \draw[gray!60, thin] (n.east) -- (-3.278,3.907);
  \node[anchor=east] (n) at (-3.60,3.30) {gap middle};
  \draw[gray!60, thin] (n.east) -- (-0.578,3.276);
  \node[anchor=west] (n) at (2.90,2.90) {target: 0.6 deepest + 0.4 middle};
  \draw[gray!60, thin] (n.west) -- (-1.080,2.031);
  \node[anchor=west] (n) at (2.90,2.20) {room ahead: 3rd percentile of ranges};
  \draw[gray!60, thin] (n.west) -- (0.235, 3.748586362432566);
  \node[anchor=west] (n) at (2.90,1.90) {inside the car's own 0.47 m corridor};
  \node[anchor=west] (n) at (2.90,0.40) {car, lidar at the origin};
  \draw[gray!60, thin] (n.west) -- (0.135, 0.0);
  \node[anchor=west] (n) at (2.90,-1.20) {beams behind the car: outside the gap search};
  \draw[gray!60, thin] (n.west) -- (0.700,-0.404);
\end{tikzpicture}
\end{center}

\begin{enumerate}[leftmargin=*]
  \item \textbf{Clean the scan.} Any beam that returned nothing (infinity) is set to the lidar's
  maximum range, and every range is clipped to the sensor's limits. This is the numpy array
  \texttt{ranges}; \texttt{angles} is the matching bearing of each beam.
  \item \textbf{Find the nearest return and blank a bubble around it.} Only beams within
  $\pm$\texttt{gap\_fov} (1.6 rad, about $\pm 92^\circ$) count. The shortest of them is the
  nearest obstacle, the red dot. Every beam whose bearing is within
  $\arctan(\texttt{bubble}/r_{\min})$ of that beam is set to zero: at $r_{\min}=1.0$~m and
  \texttt{bubble}$=0.4$~m that is $\pm 22^\circ$, the red beams. A zeroed beam can never be
  chosen as a direction, so the car never aims through a gap that passes within 0.4~m of the
  closest thing it sees. That is the whole safety idea, and it is why the car brushes walls when
  \texttt{bubble} is too small and refuses corners when it is too big.
  \item \textbf{Pick the widest gap.} The non-zero beams are split into runs of consecutive
  beams; the longest run is the gap (blue). Inside it, two candidates: the deepest beam (the
  longest range) and the middle beam of the run. The target bearing is a blend,
  \texttt{deep\_weight} $\times$ deepest $+$ $(1-\texttt{deep\_weight})$ $\times$ middle, $0.6/0.4$
  by default. Pure deepest cuts corners tight and hunts; pure middle is smooth but wide.
  \item \textbf{Turn the bearing into steering.} The car is still turning at its current
  steering angle for the actuator delay (0.25~s) plus one scan period before the new command
  takes effect, so the heading change over that time is subtracted from the target first
  (delay compensation). Then \texttt{steer\_gain} scales the bearing (0.7: a $20^\circ$ target
  asks for $14^\circ$ of steering), the result is clipped to the $\pm 30^\circ$ limit, low-passed
  with \texttt{steer\_tau} and rate-limited to the actuator's 3.2~rad/s. Without the filtering,
  a target that jumps from one beam to the next would flick the wheel full lock every scan.
  \item \textbf{Measure the room ahead.} Take the beams whose end point lies inside the car's
  own corridor (a strip 0.235~m either side of the centreline, the green dashes). The 3rd
  percentile of their forward distance is the clearance: not the minimum, which one noisy beam
  would spoil, but close to it. It is low-passed, except that a genuinely close obstacle
  (under \texttt{brake\_margin} $+0.3$~m) passes straight through so braking is immediate.
  \item \textbf{Pick a speed.} Three ceilings: \texttt{max\_speed} (2.5~m/s); the cornering
  speed $\sqrt{\texttt{lat\_accel}/\kappa}$ from the slowly filtered steering, where
  $\kappa=\tan(\delta)/L$ is the path curvature; and the braking speed
  $\sqrt{2\,\texttt{decel\_limit}\,(\mathrm{clearance}-\texttt{brake\_margin})}$, the fastest speed
  from which the car can still stop before the clearance. The goal is the smallest of the
  three, floored at \texttt{min\_speed} (0.9~m/s) so the car never parks in a corner. The goal
  is low-passed and then slewed at \texttt{accel\_limit} up and \texttt{decel\_limit} down.
  \item \textbf{Publish.} The speed target and the steering angle go out on \texttt{/drive}.
  Under the node, the supervisor clips the speed to the mode's cap and the sim adapter turns it
  into throttle: feed-forward $v/\texttt{speed\_per\_throttle}$ (23~m/s per unit throttle,
  measured with the sysid tools) plus a small PI trim clipped to half of it, slew-limited, so the
  command never drops to the simulator's hard-brake zero while speed is wanted. The same law
  used to live in every driver; now no driver needs it, and on the car the VESC does the job.
  On shutdown the node publishes speed $0$ once so the car stops.
\end{enumerate}

\textbf{Follow mode} (\texttt{follow\_gap} $> 0$, used by the house racer as the opponent):
if something sits within that distance in a $\pm 15^\circ$ cone straight ahead, it is treated
as a car to sit behind, not a wall to steer around. Its beams are replaced by the corridor seen
beside it, so the steering keeps the corridor's line, and the speed goal is capped to hold the
distance.

\textbf{Where to start tuning.} \texttt{bubble} and \texttt{deep\_weight} decide the line;
\texttt{max\_speed}, \texttt{lat\_accel} and \texttt{brake\_margin} decide how fast it dares.
Change one number, run the same command, watch the lap time. The sysid tools are for when
you want the adapter's throttle law to match the car rather than guess.

